\begin{textsample}{POS Dim 1 – 2021\_c6 – Score -1.00 – 2021\_c6/enade\_2021\_des\_64eae85373.txt}  \label{ex:f1_pos_203}
CAPÍTULO 6 ANÁLISE TÉCNICA DA PROVA Este capítulo tem por objetivo apresentar o desempenho dos estudantes concluintes de Design no Enade/2021.
Para isso, foram calculadas as estatísticas básicas da prova em seu todo ( seção 6.1.1 ), bem como as estatísticas relacionadas aos Componentes de Formação Geral ( seção 6.1.2 ) e de Conhecimento Específico da Área ( seção 6.1.3 ).
Dadas as suas características, foram analisadas, em separado, as questões objetivas ( seção 6.2 ) e as questões discursivas ( seção 6.3 ).
Tomando -se como base as duas Questões Discursivas do Componente de Formação Geral, nas seções 6.3.1.6 a 6.3.1.8, são apresentados os resultados e os comentários sobre a correção das respostas em relação à Língua Portuguesa.
Para as questões objetivas, foram disponibilizados os Índices de Facilidade e de Discriminação Ponto-Bisserial, também em separado, para os Componentes de Formação Geral ( seção 6.2.1 ) e de Conhecimento Específico ( seção 6.2.2 ).
De cada componente, uma das questões foi escolhida para exemplificar a análise gráfica, relacionando -se as alternativas escolhidas pelos estudantes ( inclusive o gabarito ) com o número de acertos no componente.
No Anexo I, apresenta -se a íntegra da análise gráfica para todas as questões objetivas.
Para cada uma das questões discursivas, os conteúdos dos tipos mais comuns de respostas dos estudantes são apresentados e comparados com o padrão de resposta esperado ( ver Anexo IX com o padrão de respostas ).
Nas tabelas, constam as seguintes estatísticas das notas¹⁹ : média do desempenho na prova, erro padrão da média, desvio padrão, nota mínima, mediana e nota máxima para cada um de seus componentes.
Tais estatísticas contemplam o total de estudantes concluintes da Área de Design inscritos e presentes à prova do Enade/2021, tendo em vista agregações, ou por Grandes Regiões e o país como um todo, ou por Categoria Administrativa, Organização Acadêmica da IES e Modalidade de Ensino.
Em relação aos gráficos de distribuição de notas, o intervalo considerado foi de 10 unidades, aberto à esquerda e fechado à direita, com exceção do primeiro intervalo, [ 0 ; 10 ], fechado em ambos os extremos.
Para os gráficos de distribuição das notas das questões discursivas, foram consideradas mais duas categorias : questão em branco e nota zero. ¹⁹ Uma definição dessas estatísticas pode ser encontrada no Glossário. 6.1 ESTATÍSTICAS BÁSICAS DA PROVA Esta seção apresenta estatísticas selecionadas e histogramas da nota geral ( 6.1.1 ) e de cada componente : Formação Geral ( 6.1.2 ) e Conhecimento Específico ( 6.1.3 ).
São também apresentadas estatísticas selecionadas de subpopulações, caracterizadas por Grande Região, Categoria Administrativa, Organização Acadêmica e Modalidade de Ensino. 6.1.1 ESTATÍSTICAS BÁSICAS GERAIS Na Tabela 6.1, são apresentadas as Estatísticas Básicas da prova, por Grande Região, dos estudantes concluintes de Design.
A Média das notas da prova, como um todo ( nas seções seguintes serão analisados os Componentes de Formação Geral e de Conhecimento Específico ), foi 48,4, sendo que os estudantes da região Norte obtiveram a Média mais baixa ( 46,4 ), e os da região Centro-Oeste obtiveram a Média mais alta ( 56,3 ).
As demais Médias foram : 47,6, na região Nordeste, 47,7, na região Sudeste e 49,6, na região Sul.
O Desvio padrão para o Brasil, como um todo, foi 12,8, sendo o maior Desvio padrão encontrado na região Nordeste ( 13,0 ), e o menor, na região Norte ( 11,6 ), indicando uma dispersão um pouco menor das notas dessa última região.
A região que obteve a maior nota Máxima foi a Nordeste ( 87,1 ), ao passo que a região que atingiu a menor nota Máxima foi a Norte ( 68,4 ).
A Mediana do Brasil, como um todo, foi 49,0, sendo a maior Mediana obtida na região Centro-Oeste ( 58,9 ), e a menor obtida na região Norte ( 47,4 ).
A nota Mínima do Brasil foi obtida nas regiões Nordeste e Sudeste, onde foi zero.
A maior nota Mínima, entre as cinco Grandes Regiões, foi 28,0, na região Centro-Oeste, as demais regiões obtiveram as seguintes notas Mínima : 10,2 ( Norte ) e 4,1 ( Sul ).
Considerando -se as notas segundo Grande Região, observa -se que existe diferença estatisticamente significativa ao nível de 95 \% entre as Médias obtidas na região Centro-Oeste e nas demais regiões, assim como entre a região Sul e as regiões Nordeste e Sudeste.
O comportamento das notas dos estudantes de todo o Brasil pode ser observado no Gráfico 6.1, que apresenta um histograma com a distribuição das mesmas : uma distribuição unimodal com a moda no intervalo ( 50 ; 60 ].
Na Tabela 6.2, são apresentadas informações referentes à Média da nota final, desagregadas por Categoria Administrativa, por Organização Acadêmica e por Modalidade de Ensino.
A Média das notas da prova, como um todo, foi 48,4.
Em relação a Categoria Administrativa, os estudantes das IES Públicas obtiveram Média mais alta ( 51,3 ) que a Média nacional, e os das IES Privadas obtiveram Média mais baixa ( 46,7 ).
Observa -se que existe diferença estatisticamente significativa entre as Médias das notas das IES Públicas e as das Privadas.
A diferença entre as Médias das regiões Centro-Oeste e Norte ( 9,9 ), a maior e a menor Média, é superior à diferença entre a das IES Públicas e Privadas ( 4,6 ), chegando a caracterizar uma maior diversidade administrativa do que regional.
O Desvio padrão para as IES Públicas ( 12,7 ) foi inferior ao do Brasil, como um todo ( 12,8 ), indicando uma dispersão um pouco menor das notas nesta Categoria Administrativa.
No tocante à Organização Acadêmica, as Universidades ( 48,7 ) obtiveram Média mais alta que a nacional ( 48,4 ).
A Média das Faculdades foi igual à nacional.
E a dos Centros Universitários ( 46,6 ) e CEFET/IF ( 42,2 ) foi menor do que a nacional.
Constata -se que existe diferença estatisticamente significativa, ao nível de 95 \%, entre as Médias das Universidades em relação aos Centros Universitários e CEFET/IF.
A Modalidade Educação a Distância não teve cursos de Design avaliados e, portanto, os valores referentes à Modalidade Educação Presencial são equivalentes ao total nacional. 6.1.2 ESTATÍSTICAS BÁSICAS NO COMPONENTE DE FORMAÇÃO GERAL Na Tabela 6.3, são apresentadas as Estatísticas Básicas em relação ao componente da prova que avalia a Formação Geral dos estudantes concluintes.
A nota de Formação Geral é obtida levando -se em conta os acertos nas oito questões objetivas e a nota média obtida nas Questões Discursivas 1 e 2, sendo a nota de cada questão discursiva obtida levando -se em conta 80 \% da nota relativa ao conteúdo e 20 \% da nota relativa ao desempenho linguístico.
Os estudantes de todo o Brasil obtiveram desempenho médio de 44,3.
Quanto à variabilidade, o Desvio padrão das notas dos estudantes do Brasil foi 15,7.
A maior Média foi obtida na região Centro-Oeste ( 53,6 ) e a menor, na região Norte ( 39,2 ).
As demais Médias foram : 43,7, na região Nordeste, 43,3, na região Sudeste e 46,1, na região Sul.
Já o maior Desvio padrão foi obtido na região Sul ( 16,1 ), e o menor, na região Norte ( 14,5 ).
Os demais desvios foram : 15,6, na região Nordeste, 15,4, na região Sudeste e 15,0, na região Centro-Oeste.
A maior nota no Componente de Formação Geral da prova do Enade/2021 foi 86,4, obtida por, pelo menos, um estudante na região Sul.
A menor nota Máxima foi obtida na região Norte ( 73,6 ).
A Mediana do Brasil, como um todo, foi 45,0, sendo a menor Mediana encontrada na região Norte ( 40,1 ), e a maior, na região Centro-Oeste ( 55,2 ).
A nota Mínima nessa parte foi zero em quase todas as regiões, exceto na região Centro-Oeste ( 13,6 ).
Considerando -se as notas, segundo Grande Região, observa -se que não existe diferença estatisticamente significativa entre as Médias das notas do Componente de Formação Geral obtidas nas regiões Nordeste e Sudeste.
As diferenças entre as demais Médias são estatisticamente significativas.
No Gráfico 6.2, é possível verificar a avaliação do desempenho dos estudantes do Componente de Formação Geral, a partir do histograma da distribuição das notas correspondentes.
A distribuição é unimodal, com moda em ( 40 ; 50 ], intervalo imediatamente anterior ao modal da distribuição de notas da prova, como um todo ( Gráfico 6.1 ).
Nota -se, ainda, que, no Gráfico 6.2, as notas apresentam maior dispersão do que as no Gráfico 6.1 ( distribuição das notas da prova ), confirmada pela comparação dos desvios padrões : 12,8 para a nota da prova, como um todo, e 15,7 para o Componente de Formação Geral.
Na Tabela 6.4, são apresentadas as informações referentes ao desempenho dos concluintes do Componente de Formação Geral, em diferentes agregações : Categoria Administrativa, Organização Acadêmica e Modalidade de Ensino.
Observa -se que existe diferença estatisticamente significativa entre as Médias dos tipos de Categoria Administrativa : a maior Média obtida por estudantes de IES Públicas ( 47,5 ) e a menor, pelos de IES Privadas ( 42,4 ).
Nota -se que, também, há diferença estatisticamente significativa entre as Médias dos Centros Universitários ( 41,7 ) em relação às Universidades ( 44,7 ) e Faculdades ( 44,6 ).
A menor Média das Organizações Acadêmicas ocorreu nos CEFET/IF ( 36,1 ).
A Modalidade Educação a Distância não teve cursos de Design avaliados e, portanto, os valores referentes à Modalidade Educação Presencial são equivalentes ao total nacional. 6.1.3 ESTATÍSTICAS BÁSICAS DO COMPONENTE DE CONHECIMENTO ESPECÍFICO Na Tabela 6.5, são apresentadas as Estatísticas Básicas referentes ao Componente de Conhecimento Específico da Área de Design.
A nota do componente de Conhecimentos Específicos leva em conta as notas da parte objetiva e a média das notas obtidas nas três questões discursivas da parte discursiva da Área.
A Média do desempenho dos estudantes do Brasil, como um todo, foi 49,7.
A maior Média foi obtida na região Centro-Oeste ( 57,2 ), e a menor, na região Norte ( 48,8 ).
As demais Médias foram : 48,9 na região Nordeste, 49,1 na região Sudeste e 50,8 na região Sul.
Quanto à variabilidade das notas, o Desvio padrão do Brasil, como um todo, foi 14,1, sendo o maior Desvio padrão observado na região Nordeste ( 14,6 ) e o menor, na região Norte ( 13,2 ).
Os demais desvios foram : 14,1 na região Sudeste, 13,8 na região Sul e 13,6 na região Centro-Oeste.
A Mediana das notas dos estudantes de todo o Brasil foi 50,8.
A maior Mediana ocorreu na região Centro-Oeste ( 59,3 ), e a menor, na região Sudeste ( 49,3 ).
As demais Medianas foram : 50,8 na região Norte, 49,8 na região Nordeste e 51,7 na região Sul.
A nota Máxima do Brasil, como um todo, foi 90,5, sendo obtida por pelo menos um estudante na região Nordeste.
As demais notas Máximas foram : 74,0 na região Norte, 87,2 na região Sudeste, 85,7 na região Sul e 80,5 na região Centro-Oeste.
A nota Mínima foi zero nas regiões Nordeste, Sudeste e Sul.
As demais notas Mínimas foram : 9,0 na região Norte e 27,7 na região Centro-Oeste.
Observa -se que existe diferença estatisticamente significativa entre a região Centro-Oeste e as demais regiões na análise das Médias das notas do Componente de Conhecimento Específico.
As diferenças entre as regiões Nordeste e Sul também são estatisticamente significativas Assim como nos Gráficos 6.1 e 6.2, no Gráfico 6.3, apresentado a seguir, é possível ser feita uma avaliação do desempenho de concluintes de Design, em relação ao Componente de Conhecimento Específico, com um histograma da distribuição das notas correspondentes.
Esta também é uma distribuição unimodal, e o grupo modal é o ( 50 ; 60 ], o mesmo do grupo modal da prova, como um todo, e acima do grupo modal para a Formação Geral.
Na Tabela 6.6, é apresentada uma comparação dos resultados em relação a Categoria Administrativa, a Organização Acadêmica e a Modalidade de Ensino, agora levando -se em conta o desempenho de estudantes do Componente de Conhecimento Específico da prova.
No que se refere a Organização Acadêmica, a maior Média foi a das Universidades ( 50,0 ), vindo a seguir a das Faculdades ( 49,6 ), a dos Centros Universitários ( 48,2 ) e, por fim, a dos CEFET/IF ( 44,2 ).
O maior Desvio padrão, e superior ao valor para o Brasil, como um todo, foi o das Faculdades ( 14,3 ).
Os demais desvios foram : 14,2 nas Universidades e 13,3 nos Centros Universitários e nos CEFET/IF.
As Universidades obtiveram a maior nota Máxima ( 90,5 ).
Os Centros Universitários obtiveram nota Máxima 86,8, as Faculdades, 82,0, e os CEFET/IF, 68,0.
As Medianas foram : 51,0 nas Universidades, 48,7 nos Centros Universitários, 49,8 nas Faculdades e 42,4 nos CEFET/IF.
A nota Mínima foi zero nas Universidades e nas demais regiões foi : Centros Universitários ( 12,8 ), Faculdades ( 8,5 ) e CEFET/IF ( 22,8 ).
Observa -se que existe diferença estatisticamente significativa, ao nível de 95 \%, no Componente de Conhecimento Específico apenas entre as notas das Universidades e Centros Universitários.
Quanto a Categoria Administrativa, observa -se um comportamento semelhante àquele da parte de Formação Geral e à prova, como um todo, ou seja, existe diferença estatisticamente significativa entre as Médias das IES Públicas ( 52,5 ) e as das IES Privadas ( 48,1 ).
Neste caso, também, a maior Média foi obtida por estudantes de IES Públicas de ensino.
A Modalidade Educação a Distância não teve cursos de Design avaliados e, portanto, os valores referentes à Modalidade Educação Presencial são equivalentes ao total nacional. 6.2 ANÁLISE DAS QUESTÕES OBJETIVAS Esta seção apresenta estatísticas selecionadas e histogramas das Questões Objetivas de Formação Geral ( 6.2.1 ) e de Conhecimento Específico ( 6.2.2 ).
São, também, apresentadas e comparadas as médias das subpopulações caracterizadas por Grande Região. 21 6.2.1 COMPONENTE DE FORMAÇÃO GERAL Na Tabela 6.7, são apresentadas as Estatísticas Básicas relativas às oito Questões Objetivas do componente da prova que abrange a Formação Geral dos estudantes.
A Média do Brasil foi 47,8.
A menor Média foi encontrada na região Norte ( 43,1 ), e a maior, na região Centro-Oeste ( 57,4 ).
As demais Médias foram : 48,3 na região Nordeste, 46,3 na região Sudeste e 49,6 na região Sul.
O Desvio padrão do Brasil foi 20,1, sendo o maior Desvio padrão encontrado na região Centro-Oeste ( 20,6 ), e o menor Desvio padrão encontrado na região Norte ( 19,2 ).
As Medianas do Brasil, como um todo, e de três regiões ( Nordeste, Sudeste e Sul ) foram iguais a 50,0.
Na região Norte ( 37,5 ) a Mediana foi menor e na região Centro- Oeste ( 62,5 ) foi maior.
A nota Máxima 100,0 foi alcançada em quase todas as Grandes Regiões, exceto na região Centro-Oeste ( 87,5 ).
As notas Mínimas ( 0,0 ) foram iguais para todas as regiões.
Na Tabela 6.8, são apresentados o Índice de Facilidade e o Índice de Discriminação ( Ponto-Bisserial ) para cada uma das questões objetivas do Componente de Formação Geral.
Quanto ao Índice de Facilidade, foram usadas as seguintes cores para diferenciar o nível de dificuldade da questão :  Azul para as questões classificadas com índice Muito fácil ( \&gt ; =0,86 ), verde para as questões classificadas com índice Fácil ( 0,61 a 0,85 ), amarelo para as questões classificadas com Médio ( 0,41 a 0,60 ), vermelho para as questões classificadas com Difícil ( 0,16 a 0,40 ) e roxo para as questões classificadas com Muito difícil ( \&lt ; =0,15 ).
Já quanto ao Índice de Discriminação, foram usadas as seguintes cores para qualificar a questão :  As questões classificadas com índice Fraco receberam a cor vermelha ( \&lt ; =0,19 ), as classificadas com Médio receberam a cor amarela ( 0,20 a 0,29 ), as classificadas com Bom receberam a cor verde ( 0,30 a 0,39 ) e as classificadas com Muito bom ( \&gt ; =0,40 ) receberam a cor azul.
As Questões Objetivas do Componente de Formação Geral, segundo o Índice de Facilidade, foram assim avaliadas : nenhuma das oito questões foi anulada pela Comissão Assessora de Área.
Nenhuma teve o Índice de Facilidade classificado como Muito fácil, e duas questões foram tidas como Fácil, por terem índice de acertos entre 0,61 e 0,85.
Outras duas questões foram consideradas com índice de dificuldade Médio, situando -se no intervalo entre 0,50 e 0,59 do Índice de Facilidade, ou seja, houve entre 50,0 \% e 59,0 \% de acertos, enquanto quatro questões foram classificadas na categoria Difícil, apresentando acertos na faixa entre 16 \% e 40 \%.
Nenhuma questão apresentou menos de 14 \% de acertos, razão pela qual não houve questão classificada como Muito difícil.
Na Tabela 6.9, é apresentada a distribuição das questões, levando -se em conta, simultaneamente, a classificação dos dois índices.
Como já comentado, nenhuma questão foi anulada.
As seis questões que tiveram Índice de Discriminação Muito bom figuraram em três níveis de dificuldade Fácil, Médio e Difícil : duas classificadas na categoria Fácil ( Questões 6 e 7 ) do Índice de Facilidade, outras duas, na categoria Médio ( Questões 1 e 3 ), e mais duas, na categoria Difícil ( Questões 4 e 5 ).
Em particular, a Questão 1 foi a que apresentou o maior poder discriminatório, com índice 0,49, e foi considerada Médio em termos de Facilidade, com uma proporção de 0,58 acertos.
O máximo de acertos foi alcançado pela Questão 7, com um Índice de Facilidade de 0,71.
A Questão 8 foi considerada Difícil, com Índice de Facilidade 0,23.
O Índice de Discriminação foi Bom para a Questão 8 ( 0,30 ).
No Gráfico 6.4, para exemplificar, analisa -se o comportamento da Questão 3 de Formação Geral.
Trata -se de uma questão considerada com nível Médio de dificuldade, em O Índice de Facilidade variou de 0,23 a 0,71, e o de Discriminação, de 0,30 a 0,49.
Seis questões, quanto ao Índice de Discriminação, foram tidas como Muito bom e duas questões foram tidas como Bom.
Na Tabela 6.9, é apresentada a distribuição das questões, levando -se em conta, simultaneamente, a classificação dos dois índices.
Como já comentado, nenhuma questão foi anulada.
As seis questões que tiveram Índice de Discriminação Muito bom figuraram em três níveis de dificuldade Fácil, Médio e Difícil : duas classificadas na categoria Fácil ( Questões 6 e 7 ) do Índice de Facilidade, outras duas, na categoria Médio ( Questões 1 e 3 ), e mais duas, na categoria Difícil ( Questões 4 e 5 ).
Em particular, a Questão 1 foi a que apresentou o maior poder discriminatório, com índice 0,49, e foi considerada Médio em termos de Facilidade, com uma proporção de 0,58 acertos.
O máximo de acertos foi alcançado pela Questão 7, com um Índice de Facilidade de 0,71.
A Questão 8 foi considerada Difícil, com Índice de Facilidade 0,23.
O Índice de Discriminação foi Bom para a Questão 8 ( 0,30 ).
No Gráfico 6.4, para exemplificar, analisa -se o comportamento da Questão 3 de Formação Geral.
Trata -se de uma questão considerada com nível Médio de dificuldade, em relação à Facilidade, e a que obteve o segundo maior Índice de Discriminação dessa parte da prova.
Neste gráfico, cada uma das cinco curvas representa o percentual de respostas em determinada alternativa da questão, em função do número de acertos dos estudantes nessa parte da prova ( Formação Geral/Múltipla Escolha ), antes de possíveis eliminações pelo critério Ponto-Bisserial.
Como já foi mencionado, nenhuma questão objetiva do Componente de Formação Geral foi anulada pela Comissão Assessora de Área, por isso o número de acertos alcança no máximo oito para esse componente.
A curva em verde corresponde à alternativa B, a alternativa correta para essa questão.
Observa -se que entre os estudantes com menor número de acertos, nessa parte do exame, a situação mais frequente foi a escolha da alternativa E ( em vermelho ), incorreta.
Por exemplo, entre os estudantes que acertaram duas respostas, 39,8 \% escolheram a alternativa E ( em vermelho ), 30,3 \% escolheram a alternativa B ( em verde, gabarito ), 13,3 \% escolheram a alternativa D ( em laranja ), 11,4 \% escolheram a alternativa A ( em azul ) e 5,2 \% escolheram a alternativa C ( em preto ).
Dentre os que acertaram duas respostas entre as questões de múltipla escolha de Formação Geral, nenhum ( 0,0 \% ) deixou a questão em branco ou marcou mais de uma alternativa, invalidando a questão.
À medida que o número de acertos aumenta, indicando desempenho melhor nessa parte da prova, aumenta, concomitantemente, a proporção de estudantes que selecionaram a alternativa correta B, atingindo 100 \% para os estudantes com oito acertos.
Essa análise permite verificar como a questão discriminou os grupos de desempenho, justificando -se o alto índice ( 0,49 ) obtido na questão.
Cumpre notar que não é possível inferir deste gráfico nem o Índice de Facilidade, que seria uma média da proporção ponderada pela quantidade de estudantes com cada uma das notas, nem o Índice de Discriminação Ponto-Bisserial, por razão equivalente.
No caso extremo, no qual a grande concentração dos acertos dos estudantes fosse abaixo de quatro, o Índice de Facilidade seria obrigatoriamente abaixo de 70,0 \% ( neste exemplo ).
Caso a concentração fosse em seis acertos ou mais, o índice seria obrigatoriamente acima de 90,0 \%.
Os gráficos relativos às demais questões de Formação Geral constam do Anexo I. 6.2.2 COMPONENTE DE CONHECIMENTO ESPECÍFICO Na Tabela 6.10, são apresentadas as Estatísticas Básicas em relação às Questões Objetivas do Componente de Conhecimento Específico da prova de Design por Grande Região.
Nessa parte da prova, como será discutido ainda nesta seção, nenhuma Questão Objetiva do Componente de Conhecimentos Específicos foi anulada pela CAA e sete questões não foram usadas no cômputo das notas por terem sido descartadas pelo critério do Ponto-Bisserial.
Assim, as notas foram calculadas com base em 20 das 27 questões objetivas válidas de Conhecimentos Específicos.
A Média do Brasil deste componente foi 53,8.
A menor Média foi observada na região Sudeste ( 53,0 ), e a maior, na região Centro-Oeste ( 62,2 ).
O Desvio padrão de todo o Brasil foi 15,7, sendo o menor Desvio padrão encontrado na região Norte ( 14,9 ), e o maior, na região Nordeste ( 16,1 ).
A Mediana de todo o Brasil foi 55,0, o mesmo valor da Mediana encontrada em quase todas as regiões, com exceção da região Centro-Oeste ( 65,0 ).
A nota Máxima da prova foi 100,00, obtida nas Questões Objetivas do Componente de Conhecimento Específico, por, pelo menos, um estudante da região Nordeste.
Nas demais regiões o valor foi : 80,0 na região Norte, a menor, 95,0 nas regiões Sudeste e Sul, e 90,0 na região Centro-Oeste.
A nota Mínima foi zero, obtida por pelo menos um estudante em quase todas as regiões, com exceção das regiões Centro-Oeste ( 25,0 ) e Norte ( 10,0 ), nas Questões Objetivas do Componente de Conhecimento Específico.
Na Tabela 6.11, são apresentados os Índices de Facilidade e de Discriminação ( Ponto-Bisserial ) das questões objetivas do Componente de Conhecimento Específico, para os estudantes de Design.
Para facilitar a diferenciação das questões, serão usadas as mesmas cores da Tabela 6.8 para as diferentes classificações dos Índices de Facilidade e de Discriminação.
Observa -se que nenhuma questão foi anulada pela CAA, como já mencionado.
Sendo assim, as análises a seguir dizem respeito às 27 questões válidas.
A partir do Índice de Facilidade obtido, pode -se concluir que pouco mais de dois quintos das questões objetivas da prova foram consideradas Difícil ou Muito Difícil : das 27 questões válidas, nove foram classificadas como Difícil e duas como Muito difícil.
Sete questões foram classificadas como Fácil, e outras nove consideradas como Médio.
Nenhuma questão foi classificada como Muito fácil.
Já quanto aos Índices de Discriminação das Questões Objetivas do Componente de Conhecimento Específico da prova, obtém -se como resultado a seguinte classificação : nove das 27 questões válidas foram consideradas boas, enquanto três delas tiveram Índice de Discriminação Muito bom.
Assim, para doze em 27 questões, os Índices de Discriminação foram Bom ou Muito bom.
Dentre as demais, oito delas foram classificadas como Médio, e outras sete, como Fraco, sendo 15, por conseguinte, a quantidade de questões nos dois patamares mais baixos de Discriminação.
Constata -se, assim, que a prova – no que se refere ao Componente de Conhecimento Específico – possuía média capacidade de discriminar entre aqueles que dominam ou não o conteúdo.
O Índice de Facilidade variou de 0,11 a 0,85 e o de Discriminação, de -0,07 a 0,47.
Na Tabela 6.12, é apresentada a distribuição das questões levando -se em conta, simultaneamente, a classificação dos dois índices.
Entre as questões que alcançaram os maiores Índices de Discriminação, a maior delas, a de número 32, foi classificada com Índice de Discriminação Muito bom.
Seu índice foi 0,47, e, quanto ao Índice de Facilidade, foi classificada como Médio, com 54 \% dos estudantes marcando a opção correta.
A Questão de número 35 foi a mais difícil entre as 27 questões específicas válidas, com baixo Índice de Facilidade, apenas 11 \% de acertos.
Essa questão apresentou poder discriminatório muito baixo, 0,03, o que comprova ter sido a mais difícil para os estudantes.
Destaca -se, também, a Questão 33, com Índice de Facilidade 0,12, o que, em termos percentuais, corresponde a 12 \% de estudantes que responderam acertadamente, sendo, também, 0,01 seu Índice de Discriminação.
Tais questões foram, portanto, pelo critério Ponto-Bisserial, consideradas inadequadas.
Por isso, as Questões 33 e 35, além das Questões 9, 15, 16, 30 e 34, foram eliminadas do cômputo da nota final.
Bisserial, consideradas inadequadas.
Por isso, as Questões 33 e 35, além das Questões 9, 15, 16, 30 e 34, foram eliminadas do cômputo da nota final.
A título de exemplo das análises do comportamento das questões objetivas, no Gráfico 6.5, analisa -se a Questão 32 do Componente de Conhecimento Específico.
Essa questão foi considerada, pelas respostas dos estudantes avaliados na prova, como uma questão com nível de Facilidade Médio, 0,54, ou seja, 54 \% dos estudantes assinalaram, acertadamente, a opção A correspondente ao gabarito.
Como já comentado, seu Índice de Discriminação foi igual a 0,47, classificado como Muito bom, o maior valor de Discriminação.
Neste gráfico, cada uma das cinco curvas representa o percentual de respostas em determinada alternativa da Questão 32, em função do número de acertos dos estudantes nessa parte da prova, antes de possíveis eliminações de questões pelo critério Ponto-Bisserial.
A alternativa correta E, representada no gráfico pela curva em vermelho, foi escolhida em maiores proporções pelos estudantes com desempenho melhor nessa parte da prova.
Já as alternativas incorretas, também denominadas distratores, foram selecionadas, principalmente, por aqueles com notas mais baixas.
No caso de Design, como se observa no eixo horizontal do Gráfico 6.5, nenhum estudante acertou mais do que 23 questões dentre as 27 questões objetivas válidas, considerando -se o gabarito original.
A eliminação das sete questões com Índice de Discriminação Fraco do cômputo da nota é que possibilitou a existência de nota Máxima igual a 100,0 por pelo menos um estudante da região Nordeste ( visto na Tabela 6.10 ).
Observa -se que a soma não é 100 \%, por causa das questões não respondidas ou com mais de uma opção marcada.
Aqueles com um acerto, na sua totalidade, marcaram a opção A ( linha azul ).
A proporção de estudantes que selecionaram a resposta correta E aumenta gradativamente, chegando a atingir 100 \% para 21 acertos ou mais, enquanto a proporção dos que escolheram alternativas incorretas decai, a partir de cinco acertos, em função do número de acertos nessa parte da prova.

% matched lemmas: 
\end{textsample}
